\clearpage
\section{Spec185：准备模型对象与应用入口}
\lead{PLANNED · DI；接口与完整任务见 specs/185-prepared-model-runtime}
应用通过 Runtime 创建 User，调用 prepare 得到 PreparedModel，再以输入和可选 placement
提交 request。NativeInferenceClient 保留唯一原生执行责任，过渡期保留旧签名。
准备成功只表示本地不可变模型包可用，不表示权限有效、Provider 已选择或远端 runner 已就绪。

\subsection{API 与身份}
Runtime::open 由 C++ 解析操作员配置并冻结模型注册；普通 User::prepare 只接模型 key，默认 default。
高级目录/source 接口单独分层。PreparedModel 持有已验证目录、splitter 与 source lease；manifest 和 receipt 只读。
request 接收 owning Input 和 RequestOptions；placement 只有 options 中一个覆盖入口。
RequestHandle 提供 result、cancel、可退订 completion 和可靠 EventReader；observe 仅作可丢弃诊断。
局部等待超时不取消请求。能力/schema 可查询，结果和错误采用独立公开类型。
完整签名、字段、默认值及异常见本 Spec 的 contracts/public-api.md。

\subsection{缓存和授权}
默认 UseOrFetch：命中返回、有同键准备则加入、否则启动准备。RequireReady 只接受 READY；
UseOrWait 仅加入已有准备，无任务时立即失败；Refresh 原子换代，失败保留旧 lease。
缓存键覆盖模型、任务、目录、adapter、splitter、source 与安全域；并发等待者独立超时和取消。
预算满且无可驱逐对象时明确失败，不留下半成品。目录复用不省略每请求输入与模型绑定检查。
动态候选、ACK admission、placement、grant、Selection 和 plan 不进入共享缓存。

\subsection{对象与调用关系}
Runtime 拥有 IO、缓存和注册 owner；子对象持有共享 State，registry 反向使用弱引用。
PreparedModel、RequestHandle 与 Conversation 分别持有必要的模型 lease。
阻塞 prepare、wait、drain 禁止在 IO 或通知回调线程调用；close 只启动关闭，drain 是本地清理屏障。
会话复用既有 coordinator、receipt/control 和 durable journal，不创建第二个成功状态。
Provider 认证 Selection 后才取得受保护 source 并组装；不可变 artifact 可复用，KV 和请求状态隔离。
不增加全局 ASSEMBLING barrier，不开放绕过认证的 prewarm 开关。

\subsection{实施与证据}
共十一批、十六任务。先安装 C++ SDK/ABI 门，再 Runtime、准备、扩展注册、请求、会话、Provider、迁移；
完整 C++ 资格通过后才包装 Python，最后同步设计。六层 API 暴露清单隔离应用、Provider、扩展、authority、内部与兼容。
逐任务只读静态门通过后继续同批，批末组合审查后统一构建和相关测试。
C++ 生产及跨进程证据先于 Python 薄封装检查；Spec184 未完成资格继续留在原任务中。
本节不更新冻结目标 API 清单为假实现，详细新增目标契约由 Spec185 独立维护。

\subsection{独立 C++ 入口}
安装 api.hpp/provider.hpp 足以编写外部应用；C++ 提供 prepareAsync、可靠 completion、
流读取与 drainAsync，退订和清理竞争由原生 owner 处理。Provider-only Runtime 不要求 User 目录。
Python 仅做参数、异常、GIL 与 asyncio 映射，不补齐 C++ 缺失能力或维护第二套状态机。
外部安装消费与独立进程覆盖全部模式，核对 ELF/子进程无 Python 依赖；完整契约见 C-05/C-06。

\subsection{完整 API 与生命周期}
C-07 统一列出全部新稳定应用/Provider API、Python 对应入口和公开值类型；旧接口全声明索引独立保留。
普通调用仅 open、prepare(key)、run(input)；核心 Python 对象直接绑定 C++，便利层不拥有领域状态。
Subscription 统一完成、观察及单次异步读的退订；流失败不能伪装 EOF，局部等待超时不取消业务。
Runtime 外壳析构启动 close，子对象持安全 State；确定性退出仍用 close/drain。
关闭后保留同步结果读取，拒绝新增业务异步订阅；已清理 owner 的 drainAsync 有明确调用线程完成快路径。
所有接口的安装消费、析构/并发反例与绑定映射均需验收；这些是 PLANNED 目标，尚未证明独立可用。

\section{TG-01：授权版本与影响范围}
\lead{PLANNED · Core；当前基线仍采用全局 ControllerVersion}
目标是让无关主体不因一次局部授权调整而清除业务状态。ServiceController::grant、revoke、
getPolicyStatus 和 User/Provider 的状态安装是迁移边界；现有签名先保持兼容。
新状态契约须分别表达密钥版本、服务策略版本、主体授权版本及受影响服务/主体，
不能把目录 epoch、请求 attempt 或通知序号当作授权版本。

\subsection{API 行为与兼容要求}
grant 成功后仅受影响主体更新授权视图；未变的权限不触发 DKEY 获取。
密钥轮换仍可影响全部使用该密钥的参与者，不能因局部通知而跳过必要撤销。
getPolicyStatus 返回带签名的权威状态；新增字段必须定义旧客户端的保守处理及不支持版本错误。
状态安装应先验证签名、主体/服务绑定、单调版本和有效期，再原子切换。
通知接口仅携带版本线索和影响范围，收到通知不直接安装授权。
NDN-SVS 可作为通知承载候选，是否采用由后续 Spec 决定；必须保留断连重连后的权威拉取。

\subsection{验收与未决项}
验证局部 grant、全局轮换、撤销后重授、乱序/重复通知、丢通知、离线恢复和旧客户端。
记录每个节点的安装版本、权限刷新次数与密钥刷新次数；“无关节点不更新”须有定向证据。
版本字段的线格式、密钥轮换粒度和兼容协商尚未定稿，本章不声称现有实现满足目标。

\subsection{版本的作用域与判定示例}
以下是逻辑字段职责，字段名为设计占位，不是已冻结的 wire schema：keyVersion 标识实际
加密材料，servicePolicyVersion 标识服务规则，subjectAuthorizationVersion 标识主体在
服务中的授权视图，affectedSubjects/affectedServices 指明影响范围。每个比较都必须先
确认 Controller、主体和服务归属；不得跨域比较整数。

例：只给 P1 增加 S 权限时，P2 收到提示仍可验证权威状态；若 P2/S 的授权视图和密钥身份
均未变化，应保留其在途调用和已验证缓存。若密钥轮换，则使用旧密钥的 P2 仍必须刷新。
旧客户端无法解释细粒度字段时继续采用全局失效的保守路径；不能把“无法解析影响范围”
当成“不受影响”。在途请求继续还是终止须由服务安全语义明确规定。

\section{TG-02：DI 原生请求链与唯一状态所有者}
\lead{PLANNED · DI；关联 Spec182，不能据此关闭原有任务}
NativeInferenceClient 的请求入口负责驱动准备、ACK 收集关闭、候选选择、发布后绑定、
封存、提交和终态；适配器提供模型检查和拆分能力，Core 继续处理通用通信与权限。
Python APPClient 逐步调用同一原生状态机，不另行维护第二套候选和终态规则。

\subsection{独立授权进程与原生优先验收（2026-09-10）}
\lead{PLANNED · Spec182 R11-B1 至 R11-B9；本次不改变当前实现快照}
artifact authority 独立运行，持有授权签名私钥并执行 grant 策略。
requester 仅持有自身凭据、authority 公钥或信任锚及连接配置，不能读取、注入、
继承 authority 私钥，也不能在失败时退回本地 issuer。进程分离须配合文件权限
和继承句柄检查；仅有不同 PID 不构成私钥隔离。
请求经签名申请与认证答复获得 grant，绑定 request、attempt、接收者、epoch、
摘要、用途与有效期；拒绝或超时必须关闭本次授权路径，禁止继续可执行 Selection。
具体传输适配器、名字、配置、超时、取消与资源生命周期由 R11-B1 编码前冻结，
复用 Core 通信能力，不创建第二套协作协议。

执行顺序为独立 authority，随后由 \code{DI\_NativeRequester} 与
\code{di-native-provider} 完成真实跨进程 unary，观测 ACK、Selection、Provider
执行及最终 Response；再在同一原生链验证 stream、continuation、recovery、
replacement 和 cleanup。通过后才迁移原有 16 个调用方，最后执行 no-Python、
依赖闭包及 T016 整体资格。逐任务只读静态门、批次组合审查后，依次进行
C++ production code 对应的 C++ unit/integration/process tests 和 Python wrapper checks。
Python 检查数量、CLI smoke 和同进程 fixture 不替代跨进程行为证据。
阶段契约与验收反例见 Spec182 的 \code{contracts/native-first-execution.md}；
既有 T001--T017 状态和历史结果不因本次重排而改变。

\subsection{API 迁移契约}
\code{NativeInferenceHandle::result} 保留局部等待超时与请求 deadline 的区别；cancel 幂等，
observe 保留迟注册历史回放和串行通知规则。接通 dispatch 之前保留明确的
\code{NATIVE_REQUEST_PIPELINE_NOT_READY} 错误，禁止静默改走测试执行器。
\code{NativeModelSplitStrategy::enumerate} 的输出由共享候选 validator 校验；模型适配器不能
越过角色覆盖、依赖、资源和身份约束。Python 入口签名、别名、默认值与异常映射逐项登记。
\code{PlanningGraphDigest}、\code{CanonicalGraphDigest}、\code{ArtifactRootDigest} 和 RequestAttempt
建议采用不同类型，转换必须经过对应 owner 验证，禁止仅因底层都是字符串而互换。

\subsection{验收与迁移边界}
先以独立 SDK oracle 校验候选和角色映射，再检查真实 requester 调用路径、授权拒绝、
晚到 ACK、发布失败和取消竞争。组件测试、网络集成、模型执行及硬件资格分别记录。
新旧路径切换必须显式配置和可观测，兼容入口删除须先证明维护调用方已经迁移。

\subsection{对象连接与状态迁移}
\begin{tabularx}{\linewidth}{@{}p{45mm}X@{}}
\toprule
现有责任 & 目标 owner 与输出 \\
\midrule
Python APPClient / coordinator & NativeInferenceClient 持有一次 operation；Python 仅转参数、异常和事件。\\
模型/输入准备 & NativeRequestPreparation 调用固定 adapter，产出输入、真实图和候选。\\
ACK 后规划 & 原生操作接收 Core ACK\_CLOSED，admission 后调用 split/placement，并验证 proposal。\\
制品与封存 & ensureArtifacts → bindPublishedRoles → sealCore → grant/security → project/encode。\\
执行与结束 & Core 提交选择，Provider 原生 session 执行；唯一 operation 接受结果或错误，阻断旧 attempt。\\
\bottomrule
\end{tabularx}
操作应明确经历准备输入、等待 ACK、规划、确保制品、封存、提交、等待结果和终态；
在任一耗时边界检查同一 deadline/cancel，不以新子步骤重新获得完整请求预算。
Python 不再自行重建候选摘要、重试 attempt 或持有另一套终态仲裁。

\subsection{生成与会话的已发现缺口}
AC-18 记录当前 abortTurn 不持久化且未实际设置 aborted、journal 缺显式写失败/fsync
检查、beginTurn 需要父记录而缺独立首次对话引导流程。目标须为首次建会话、父记录验证、
本轮完成、Provider KV 提升、journal 提交和失败回退指定唯一 owner 与先后次序。
写失败不得返回已提交；取消必须阻止同一 attempt 的后续 checkpoint；本地恢复成功但
Provider 缓存丢失时明确重算或拒绝，不能无证据宣称 KV 命中。具体事务 API 与崩溃一致性
协议随 Spec182 的 generation/conversation 契约定稿，本轮未实现它们。

\section{TG-03：跨模块异步 API 契约}
\lead{PLANNED · Core / DI / Repo / UAV}
统一描述规则，不强制所有组件使用同一个线程池。每个异步入口登记调用线程、回调执行器、
重入许可、资源所有者、deadline 来源、取消传播范围和终态后的回调策略。
RequestServiceStreaming、NativeInferenceHandle、artifact transfer 和 UAV mission/job
优先应用此规则；实际签名及 wire 变化分别由后续 Spec 审查。

\subsection{状态和错误}
局部等待超时不隐式取消业务请求。请求 deadline 到期禁止产生新的下游工作；已经提交的
外部副作用通过显式取消/补偿处理，不能把本地句柄结束等同于远端工作已停止。
终态只能提交一次，迟到回调按 request/attempt/generation 过滤；重复取消无新增副作用。
队列达到容量时明确拒绝、丢弃旧帧或等待的策略，禁止无界增长。

\subsection{验收}
使用可控时钟和调度器覆盖 cancel/complete/close 竞争、回调重入、观察者异常、
排队超时和资源释放。远端取消须有实际协议路径证据，不能只依赖本地状态测试。

\subsection{按 API 指定策略，避免错误统一}
\begin{tabularx}{\linewidth}{@{}p{44mm}X@{}}
\toprule
入口 & 目标明确的规则 \\
\midrule
NativeInferenceHandle & 保留局部等待与请求期限区分；observer 的回放/新事件串行，cancel 幂等。\\
RequestServiceStreaming & 事件队列有界，发布拒绝原因可观察；终态一次，重传不重做业务。\\
artifact session & transfer 与 commit 分阶段取消；已提交副作用不能因本地取消自动撤回。\\
UAV mission/job & 取消任务与发送车辆停止命令分开确认；迟到旧 attempt 报告不推进新状态。\\
视频最新帧队列 & 保留丢旧帧策略，报告 dropped；不要求与完整制品同样的全量交付。\\
\bottomrule
\end{tabularx}
每项还需在后续 Spec 固定 callback executor、重入许可与 close 后是否排空队列；目前
没有全模块通用的线程安全保证。Qwen 内部状态机终态 cancel 会抛异常，兼容调整必须显式
决定是保留内部严格方法，还是由外层幂等包装；不能悄悄修改所有同名 cancel。

\section{TG-04：Repo 能力、提交与恢复}
\lead{PLANNED · Repo；保留接收、验证、提交的阶段边界}
C++ RepoCore::catalogLookup 继续表示精确对象名查询；组合过滤由单独的查询契约承载，
不改变旧接口参数含义。建立 local C++、network service、Python orchestration 三列能力表，
每项查询、上传、下载、复制和 repair 标明 supported / unsupported / partial。

\subsection{API 行为}
begin\_upload、publish\_file、begin\_fetch、fetch\_file 和 abort 应明确 operation identity、
恢复凭据、重复调用结果和进度保留策略。相同 operation 与相同摘要重试返回相同已提交结果；
相同 operation 携带不同摘要必须拒绝。目录可见性与 COMMITTED 状态关联，目录声明不能
替代实际 Data 完整性验证。条件更新和恢复记录的落盘原子性需要独立设计。

\subsection{验收}
覆盖传输中止、验证后崩溃、提交后响应丢失、重复提交、目录存在但数据缺失、删除与读取竞争。
不得把 Python repair 编排自动归入 C++ 对等能力；存储格式变更需提供迁移或版本拒绝路径。

\subsection{操作身份与条件提交}
目标操作记录应绑定逻辑操作键、制品内容摘要、manifest 身份、目标/副本集合、阶段及
验证 receipt。提交需携带预期旧状态或 revision：相同内容的已提交重试回放结果，
不同内容冲突拒绝，旧阶段请求不得覆盖新 COMMITTED/tombstone。实际字段名、数据库事务
范围与 wire 版本由后续 Spec 决定，不能以一张 DTO 表替代落盘设计。

目标能力表以三类真实入口为列：local C++、network service、Python orchestration；行至少
包括精确对象、原 packet、上传提交、恢复下载、精确目录、过滤查询、replication、repair。
当前 C++ 精确 lookup 保留支持，过滤与跨 Repo 编排不能自动升级为 C++ supported。
格式迁移先识别旧版本并验证恢复信息，再原子提交新记录；不能遇到未知版本就清空进度。

\section{TG-05：UAV 类型边界与独立状态机}
\lead{PLANNED · UAV；飞控、任务、视频与界面通过显式接口协作}
FlightControllerBackend::command/sendMavlink 的应用参数逐步由字符串字段映射改为
经验证的命令/结果类型；在兼容适配器中完成旧 Fields 转换，未知字段、缺字段、单位不符
及不支持的 schema version 明确报错。原始 MAVLink 帧通道保持独立能力和权限检查。

\subsection{API 行为}
mission 保持长期状态，有限 job 保持独立 attempt 与终态 owner。
compensateMissingParts 仅对缺失部分发起新 attempt，旧结果不得覆盖新状态。
操作租约、Core 授权和飞控 readiness 分别验证，不因 GUI 按钮可用就跳过后端检查。
视频队列的丢帧策略与任务结果交付分离；界面订阅状态，领域层不依赖 Qt 生命周期。

\subsection{验收与落地顺序}
先验证 Mock backend 的类型转换、租约过期、失联和迟到结果，再验证真实飞控适配器。
视频性能和真实飞行安全边界另行验收。本轮只接受改进方向，不宣称硬件链路已通过。
优先完成文档基线和关键契约，再沿当前 Spec 落地 TG-02；TG-01/03/04/05 各自通过
后续 Spec 确定具体接口、兼容迁移和任务，不创建虚构的 Spec 编号。

\subsection{命令模型与适配层}
候选命令模型应显式包含 schemaVersion、commandId、目标 Drone、命令类型、带单位的
参数、deadline 和 operator lease 引用；候选结果包含接受/执行/拒绝阶段、原始飞控
ACK、错误域和关联 commandId。这里定义职责，尚不固定 JSON 字段拼写与 MAVLink 映射。
旧 Fields 由兼容适配器解析为领域类型，领域层校验后才调用 backend；未知命令拒绝，
不能通过一个默认分支发送任意帧。原始帧能力继续单独授权。

mission owner 管理分片与补偿，job owner 管理有限 attempt，backend 管理设备链路与
ACK，Qt 仅订阅不可变状态并发起命令。关闭窗口不隐式宣称取消车辆动作；遥测中断时
领域层更新 readiness，结果通过明确状态传回 UI。分阶段验收先覆盖字段/单位转换和
Mock 状态竞争，再验证 UDP/串口与 SITL，最后单独评估真实飞控；不同阶段不共用一个 PASS。
